\section{Conclusions and Perspectives}\label{section:conclusion}

\subsection{Main findings}

The four parameter sweeps, the $p=1$ landscape map, and the classical baselines on the canonical $n=16$, $K=4$ instance support six concrete observations.

\begin{enumerate}
    \item Brute-force enumeration, greedy-Sharpe, rounded Markowitz, and the simulated-annealing median all recover the same four-asset basket $\{\text{GOOGL}, \text{IBM}, \text{NVDA}, \text{JPM}\}$ at Sharpe $0.136$ and $C(\boldsymbol x^\star)=-1.40\times 10^{-3}$, in milliseconds and to numerical precision (\cref{tab:classical-baselines}). The cardinality-constrained portfolio is a textbook NP-hard problem, but at $n=16$ it is not a hard \emph{instance}.
    \item QAOA at $p=3$ scores a scaled ratio $r \approx 0.94$ that is largely a ``vanity ratio'': the denominator $E_{\text{worst}} - E_{\text{opt}}$ is inflated by the budget penalty $AK^2$, while the absolute energy sits at $2847\times$ the optimum (\cref{fig:single-run}).
    \item The five most probable bitstrings at $p=3$ are all infeasible: rank 1 is the all-ones state $\mathtt{1111111111111111}$ at $0.7\%$, ranks 2--5 are the single quantum-stock states (\cref{tab:top5}). The brute-force optimum sits at rank $\sim 125$ with $p_{\text{opt}} = 1.6\times 10^{-4}$, about $10\times$ the uniform baseline $1/2^{16}$ but two decades below the leading infeasible mass.
    \item Solution quality is essentially flat across three orders of magnitude in the risk-aversion parameter $\lambda$; risk aversion is not the bottleneck on this instance.
    \item The depth sweep is non-monotonic --- $r$ peaks at $0.993$ for $p=2$, falls to $0.925$ at $p=4$, and partially recovers at $p=5$ (\cref{fig:depth-sweep}) --- consistent with the $(2\pi)^{2p}$ angle-space volume outgrowing the fixed ten-restart COBYLA budget rather than a defect in the ansatz, since the $p=1$ landscape itself is benign and basin-structured (\cref{fig:landscape}).
    \item The median top-1 success probability $p_{\text{opt}}$ collapses geometrically with $n$, from $0.23$ at $n=4$ to $1.6\times 10^{-4}$ at $n=16$, tracking $1/\binom{n}{K}$ almost as if QAOA were sampling uniformly over the feasible Hamming-weight-$K$ subspace.
\end{enumerate}

\subsection{Method-level discussion}

The exercise has clear methodological merits. Statevector simulation produces a clean, noise-free picture of QAOA dynamics suited to pedagogy, and the modular pipeline $(\boldsymbol\mu, \boldsymbol\Sigma) \rightarrow Q \rightarrow (h, J, c) \rightarrow H_C$ is dataset-agnostic. The variational principle $E \geq E_0$ supplies a built-in correctness anchor, and brute-force enumeration at $n \leq 16$ gives unambiguous ground truth against which approximation ratios can be reported honestly rather than relative to a moving target \parencite{farhi2014qaoa}.

The drawbacks are equally honest. At $n \leq 16$ the problem is classically trivial: brute force terminates in roughly $30~\text{ms}$ at $n=16$, while greedy-Sharpe and Markowitz-round reach the optimum at machine precision in under a millisecond --- about $4000\times$ faster than QAOA. The scaled ratio $r$ is a vanity number: a value of $0.94$ on a spectrum dominated by the penalty $AK^2$ coexists with an absolute energy nearly three thousand times the optimum, and with the algorithm's probability mass piling onto infeasible bitstrings (\cref{tab:top5}). The penalty mixer is also sensitive to $A$ in a way that is difficult to tune without enumerating $C(\boldsymbol x)$ --- precisely the work QAOA is meant to avoid. Ten random restarts are borderline at $p=3$ and clearly insufficient at higher depth, and the standard QAOA literature \parencite{zhou2018qaoa} prescribes warm-started angles that were not used here.

The honest framing is therefore that QAOA at this scale is valuable for the questions it lets us ask --- what $E(\boldsymbol\gamma, \boldsymbol\beta)$ looks like, how $r$ scales with $p$, and where adiabaticity breaks down --- rather than for answers it provides over Gurobi on 16-asset portfolios. No quantum advantage is claimed.

\subsection{Future directions}

Several concrete next steps would directly address the limitations identified above, forming a natural sequence from immediate fixes to longer-term comparisons. The most immediate is warm-started angle initialisation near the adiabatic schedule rather than from random points \parencite{zhou2018qaoa}, which should restore monotonicity in the depth sweep. Pairing this with fifty or more random seeds per depth would disentangle landscape difficulty from optimiser variance and give the $A$- and $\lambda$-sweeps a stable signal-to-noise ratio. Replacing the penalty mixer with a constraint-preserving XY-mixer and a Dicke initial state \parencite{mancilla2026constrained} would eliminate $A$ entirely by restricting the search to the Hamming-weight-$K$ subspace, removing the vanity-ratio artefact at its source. Extending to inequality budgets $\sum_i x_i \leq K_{\max}$ via slack variables \parencite{thomassin2025slack} is a natural follow-on once the equality case behaves well. Scaling past the $2^n$ statevector barrier looks feasible through circuit cutting \parencite{soloviev2025cutting}, which would open up problem sizes where the classical baselines are no longer trivial; benchmarking against quantum annealing on identical instances \parencite{morapakula2025annealing} would situate the results in a broader quantum landscape. A direct comparison against Gurobi or CPLEX MIP on the same instances would, finally, quantify the QAOA-vs-classical gap rather than merely assert it.
